\documentclass[11pt]{article}
\usepackage[french]{babel}
\newcommand{\imgdir}{images/fr}
% Shared layout for the Illuminati manual. Each manual-<lang>.tex loads this
% after \documentclass and sets \manuallang (babel name) and \imgdir.
% Build with XeLaTeX (see build.ps1); the fonts are the app's own.

\usepackage{fontspec}
\setmainfont{Literata}[
  Path=../../assets/fonts/literata/, Extension=.ttf,
  UprightFont=*-Regular, BoldFont=*-Bold, ItalicFont=*-Italic, BoldItalicFont=*-BoldItalic]
\setsansfont{AtkinsonHyperlegible}[
  Path=../../assets/fonts/atkinson/, Extension=.ttf,
  UprightFont=*-Regular, BoldFont=*-Bold, ItalicFont=*-Italic, BoldItalicFont=*-BoldItalic]
\newfontfamily\dyslexicfont{OpenDyslexic}[
  Path=../../assets/fonts/opendyslexic/, Extension=.otf,
  UprightFont=*-Regular, BoldFont=*-Bold, ItalicFont=*-Italic, BoldItalicFont=*-Bold-Italic]
\newfontfamily\garamondfont{EBGaramond}[
  Path=../../assets/fonts/ebgaramond/, Extension=.ttf,
  UprightFont=*-Regular, BoldFont=*-Bold, ItalicFont=*-Italic, BoldItalicFont=*-BoldItalic]

\usepackage[a4paper, margin=2.4cm, headheight=14pt]{geometry}
\usepackage{graphicx}
\usepackage[table]{xcolor}
\usepackage{booktabs, tabularx, array}
\usepackage{enumitem}
\usepackage{fancyhdr}
\usepackage{titlesec}
\usepackage[most]{tcolorbox}
\usepackage[hidelinks]{hyperref}

% Colours from the app's logo.
\definecolor{violet}{HTML}{34206A}
\definecolor{violetlight}{HTML}{EDE8F7}
\definecolor{gold}{HTML}{C9A227}
\definecolor{goldlight}{HTML}{FBF4DF}

\setlength{\parskip}{0.55em}
\setlength{\parindent}{0pt}
\setlist{itemsep=0.25em, topsep=0.3em}
\linespread{1.08}

\titleformat{\section}{\sffamily\Large\bfseries\color{violet}}{\thesection}{0.8em}{}
\titleformat{\subsection}{\sffamily\large\bfseries\color{violet}}{\thesubsection}{0.7em}{}
\titlespacing*{\section}{0pt}{1.6em}{0.6em}
\titlespacing*{\subsection}{0pt}{1.1em}{0.35em}
% Don't leave a heading alone at the bottom of a page.
\usepackage{needspace, etoolbox}
\pretocmd{\section}{\needspace{8\baselineskip}}{}{}
\pretocmd{\subsection}{\needspace{5\baselineskip}}{}{}

\pagestyle{fancy}
\fancyhf{}
\fancyhead[L]{\sffamily\small\color{violet} Illuminati}
\fancyhead[R]{\sffamily\small\color{gray}\nouppercase{\rightmark}}
\fancyfoot[C]{\sffamily\small\thepage}
\renewcommand{\headrulewidth}{0.3pt}
% \rightmark is the first section mark on the page.
\renewcommand{\sectionmark}[1]{\markboth{#1}{#1}}
\renewcommand{\subsectionmark}[1]{}

% Keep catalogue names whole in every language.
\babelhyphenation{Wiki-source Guten-berg}

% A key or key combination, e.g. \kbd{Ctrl+F}.
\newcommand{\kbd}[1]{\tikz[baseline=(k.base)]\node[draw=gray!60, fill=gray!8,
  rounded corners=2pt, inner sep=2pt, font=\sffamily\small](k){#1};}
\usepackage{tikz}

% Arrow keys (the sans font has no arrow characters).
\newcommand{\kright}{\kbd{\ensuremath{\rightarrow}}}
\newcommand{\kleft}{\kbd{\ensuremath{\leftarrow}}}
\newcommand{\kup}{\kbd{\ensuremath{\uparrow}}}
\newcommand{\kdown}{\kbd{\ensuremath{\downarrow}}}

% A button, menu item or label as it appears in the app.
\newcommand{\ui}[1]{\textsf{\textbf{#1}}}

% Screenshot with caption; skipped if the image file isn't there.
\newcommand{\shot}[2]{%
  \IfFileExists{\imgdir/#1.jpg}{%
    \begin{figure}[htbp]\centering
      \fcolorbox{gray!40}{white}{\includegraphics[width=0.96\textwidth]{\imgdir/#1.jpg}}
      \caption{#2}
    \end{figure}}{}}

\newtcolorbox{tip}[1]{enhanced, colback=goldlight, colframe=gold,
  boxrule=0.6pt, arc=3pt, left=8pt, right=8pt, top=5pt, bottom=5pt,
  fonttitle=\sffamily\bfseries, coltitle=black, colbacktitle=gold!35,
  title={#1}}
\newtcolorbox{note}[1]{enhanced, colback=violetlight, colframe=violet,
  boxrule=0.6pt, arc=3pt, left=8pt, right=8pt, top=5pt, bottom=5pt,
  fonttitle=\sffamily\bfseries, coltitle=white, colbacktitle=violet,
  title={#1}}

% Two-column tables that fill the text width.
\newenvironment{keytable}[2]{%
  \renewcommand{\arraystretch}{1.3}%
  \rowcolors{2}{white}{violetlight}%
  \tabularx{\textwidth}{@{}>{\raggedright\arraybackslash}p{0.33\textwidth}X@{}}
  \toprule \textsf{\textbf{#1}} & \textsf{\textbf{#2}} \\ \midrule}{%
  \bottomrule\endtabularx}

% Title page: \maketitlepage{subtitle}{version line}{blurb}
\newcommand{\maketitlepage}[3]{%
  \begin{titlepage}
    \centering
    \vspace*{3cm}
    \includegraphics[width=4.2cm]{../../assets/icon/app_icon.png}\par\vspace{1.2cm}
    {\sffamily\fontsize{40}{46}\selectfont\bfseries\color{violet} Illuminati\par}
    \vspace{0.5cm}
    {\sffamily\LARGE #1\par}
    \vspace{1.5cm}
    {\large\itshape #3\par}
    \vfill
    {\sffamily\color{gray} #2\par}
  \end{titlepage}}

\hypersetup{pdftitle={Illuminati — Manuel d’utilisation}, pdfauthor={Illuminati}, pdflang={fr}}

\begin{document}

\maketitlepage{Manuel d’utilisation}{Version 1.0 · Windows, macOS, Linux et Android}%
  {Trouvez, téléchargez et lisez des livres qui n’ont jamais été protégés\\par le droit d’auteur, ou dont les droits ont expiré.}

{\setcounter{tocdepth}{1}\tableofcontents}
\newpage

% ---------------------------------------------------------------------------
\section{Bienvenue}

Illuminati est une application de lecture pour Windows, macOS, Linux et Android. Elle fonctionne
comme les liseuses que vous connaissez peut-être déjà, à une différence près : tous les livres
qu’elle propose sont gratuits et légaux à télécharger, parce que leurs droits d’auteur ont expiré
ou n’ont jamais existé.

\begin{itemize}
  \item Plus de 75\,000 livres de \textbf{Project Gutenberg}, ainsi que \textbf{Standard Ebooks},
        \textbf{Wikisource} et \textbf{Internet Archive}.
  \item Un lecteur confortable avec surlignages, notes, signets, recherche, dictionnaire et
        traduction.
  \item Votre propre bibliothèque, avec des étagères et des listes de lecture.
  \item Des réglages d’affichage et d’accessibilité, dont une police conçue pour la dyslexie.
  \item Pas de compte, pas d’abonnement, pas de publicité. Tout reste sur votre appareil.
\end{itemize}

L’interface est disponible en français, anglais, espagnol, allemand, italien et portugais.

% ---------------------------------------------------------------------------
\section{Premiers pas}

\subsection{Installation}
\begin{itemize}
  \item \textbf{Windows :} copiez tout le dossier \texttt{Illuminati} sur votre ordinateur et
        lancez \texttt{illuminati.exe}. Laissez les autres fichiers à côté ; le programme en a
        besoin.
  \item \textbf{Android :} installez le fichier \texttt{.apk}. Il faudra peut-être autoriser
        l’installation d’applications de cette source dans les réglages du téléphone.
  \item \textbf{macOS et Linux :} installez la version prévue pour votre système.
\end{itemize}

\subsection{L’écran d’accueil}
Sur ordinateur, l’application s’ouvre en plein écran. L’écran d’accueil comporte deux parties :

\begin{enumerate}
  \item \ui{Votre bibliothèque} (en haut) : les livres que vous avez téléchargés, en commençant
        par le dernier lu, avec une barre de progression. Le bouton \ui{Bibliothèque et étagères}
        ouvre toute votre bibliothèque.
  \item \ui{Découvrir des livres gratuits} (en bas) : un champ de recherche, la langue du
        catalogue, les catégories, une rangée pour chaque catalogue supplémentaire et, enfin, le
        catalogue de Project Gutenberg.
\end{enumerate}

\shot{home}{L’écran d’accueil.}

Les deux boutons en haut à droite ouvrent \ui{Affichage et accessibilité}
(section~\ref{sec:display}) et \ui{À propos}, qui indique d’où viennent les livres et les licences
des polices.

% ---------------------------------------------------------------------------
\section{Trouver des livres}

\subsection{Recherche}
Tapez dans \ui{Rechercher par titre ou auteur}. La recherche :
\begin{itemize}
  \item trouve n’importe quelle partie d’un titre ou d’un nom d’auteur : \emph{misérab} trouve
        \emph{Les Misérables} ;
  \item ignore les accents et les majuscules : \emph{miserable} et \emph{MISÉRABLE} donnent les
        mêmes résultats ;
  \item interroge tous les catalogues à la fois ; chacun affiche ses résultats dans sa propre
        rangée.
\end{itemize}
Videz le champ avec le bouton \ui{×} pour revenir aux livres populaires.

\subsection{Catégories et langue}
Les étiquettes sous le champ de recherche (\ui{Fiction}, \ui{Policier}, \ui{Science-fiction},
\ui{Poésie}…) limitent tous les catalogues à un thème. Le menu de langue à droite n’affiche que les
livres d’une langue, ou \ui{Toutes les langues}. Les catalogues qui n’ont rien pour la langue ou la
catégorie choisie sont masqués.

\subsection{Les catalogues}
\begin{keytable}{Catalogue}{Ce qu’il propose}
  Project Gutenberg & La plus grande bibliothèque gratuite, avec plus de 75\,000 livres classés
                      par popularité. Beaucoup existent en édition illustrée. \\
  Standard Ebooks   & Des classiques soigneusement relus et mis en page, en anglais. \\
  Wikisource        & Des textes transcrits par des bénévoles, dans de nombreuses langues. Le
                      livre numérique est préparé au moment du téléchargement ; les œuvres très
                      longues (par exemple des romans en plusieurs tomes) peuvent être
                      indisponibles. \\
  Internet Archive  & Des numérisations d’anciennes éditions imprimées, publiées il y a 95 ans
                      ou plus. Le texte a été extrait automatiquement des numérisations et peut
                      contenir des erreurs. \\
\end{keytable}

\subsection{La page d’un livre}
Touchez une couverture pour ouvrir la page du livre. Elle affiche le résumé, les thèmes, la
langue, le catalogue et une note sur les droits d’auteur. De là, vous pouvez :
\begin{itemize}
  \item \ui{Télécharger gratuitement} : le livre est enregistré sur votre appareil, prêt à être lu
        hors connexion. Chaque couverture du catalogue a aussi un petit bouton de téléchargement.
  \item \ui{Lire maintenant} ou \ui{Reprendre la lecture}, une fois le livre téléchargé.
  \item \ui{À lire}, \ui{Terminé} et \ui{Ajouter à une étagère}, même avant de le télécharger
        (section~\ref{sec:library}).
  \item \ui{Obtenir l’édition illustrée}, pour les livres de Gutenberg téléchargés sans images.
\end{itemize}

\shot{details}{La page d’un livre.}

% ---------------------------------------------------------------------------
\section{Votre bibliothèque}\label{sec:library}

\ui{Bibliothèque et étagères}, sur l’écran d’accueil, ouvre tous les livres que vous avez
enregistrés : ceux qui sont téléchargés et ceux que vous avez gardés pour plus tard sans les
télécharger (couverture estompée et petit nuage).

\shot{library}{La bibliothèque, avec ses vues et ses étagères.}

\subsection{Vues, filtre et tri}
Les étiquettes du haut affichent une vue, chacune avec son nombre de livres : \ui{Tous},
\ui{Téléchargés}, \ui{À lire}, \ui{En cours}, \ui{Terminé}, puis chacune de vos étagères. Le champ
au-dessus filtre par titre ou auteur. Le bouton de tri (en haut à droite) classe les livres par
\ui{Lus récemment}, \ui{Titre}, \ui{Auteur} (par nom de famille), \ui{Progression} ou
\ui{Ajoutés récemment}. L’application retient votre choix.

\subsection{Listes de lecture}
\begin{itemize}
  \item \ui{À lire} garde un livre pour plus tard, sans le télécharger. Quand vous ouvrez le livre,
        il quitte cette liste.
  \item \ui{Terminé} est coché automatiquement quand vous arrivez à la fin d’un livre (vous pouvez
        annuler), ou à la main.
\end{itemize}

\subsection{Étagères}
Les étagères sont vos propres collections, comme \emph{Favoris} ou \emph{Philosophie}. Un livre
peut être sur plusieurs étagères.
\begin{itemize}
  \item Créez-en une avec \ui{Nouvelle étagère} dans la bibliothèque, ou depuis
        \ui{Ajouter à une étagère} sur n’importe quel livre.
  \item Faites un clic droit (ou un appui long) sur une étagère pour la renommer ou la supprimer.
        Quand vous supprimez une étagère, ses livres restent dans votre bibliothèque.
\end{itemize}

\subsection{Le menu d’un livre}
Faites un clic droit sur un livre (ou un appui long sur téléphone) pour voir
\ui{Détails du livre}, \ui{Ajouter à une étagère}, \ui{À lire}, \ui{Marquer comme terminé} et deux
façons de le retirer :
\begin{itemize}
  \item \ui{Supprimer le téléchargement} libère de l’espace mais garde le livre dans votre
        bibliothèque, avec ses notes, ses étagères et votre position de lecture.
  \item \ui{Retirer de la bibliothèque} supprime tout, y compris les notes et les signets.
\end{itemize}

% ---------------------------------------------------------------------------
\section{Lire}

Touchez un livre téléchargé pour l’ouvrir. Il s’ouvre là où vous l’aviez laissé.

\shot{reader}{Le lecteur et sa barre d’outils.}

\subsection{Tourner les pages}
\begin{keytable}{Action}{Comment}
  Page suivante      & Touchez le bord droit · \kright{} · \kbd{Pg Suiv} · \kbd{Espace} \\
  Page précédente    & Touchez le bord gauche · \kleft{} · \kbd{Pg Préc} \\
  Défiler un peu     & \kup{} \kdown{} ou la molette de la souris \\
  Afficher la barre  & Touchez le centre de la page \\
\end{keytable}

\subsection{La barre d’outils}
Touchez le centre de la page pour l’afficher. En haut : retour à la bibliothèque, recherche
(section~\ref{sec:search}), signet, sommaire et réglages d’affichage. En bas : chapitre précédent
et suivant, un curseur pour passer d’un chapitre à l’autre et votre progression.

\subsection{Progression et temps restant}
La ligne du bas indique combien de minutes il reste dans le chapitre et où vous en êtes dans le
livre. Le temps est calculé d’après votre propre vitesse de lecture, que l’application apprend
pendant que vous lisez.

\subsection{Sommaire et images}
\ui{Sommaire} affiche les vrais chapitres du livre, même quand un seul fichier en contient
plusieurs. Les images apparaissent dans le texte ; touchez-les pour les voir en plein écran et les
agrandir.

% ---------------------------------------------------------------------------
\section{Surlignages, notes et signets}

\subsection{Sélectionner du texte}
Sur téléphone, appuyez longuement sur un mot et faites glisser les poignées. Sur ordinateur,
faites glisser la souris sur le texte ou double-cliquez sur un mot, puis faites un
\textbf{clic droit} sur la sélection. Un menu apparaît avec :
\begin{itemize}
  \item \ui{Surligner} : marque le texte avec votre couleur de surlignage.
  \item \ui{Note} : un surlignage avec votre propre note, au choix parmi quatre couleurs.
  \item \ui{Définir} et \ui{Traduire} (section~\ref{sec:lookup}).
  \item \ui{Copier}.
\end{itemize}
Les surlignages s’ajustent aux mots entiers. Ceux qui ont une note sont soulignés. Sélectionnez de
nouveau un passage surligné pour modifier sa note ou le supprimer.

\subsection{Signets}
Utilisez le bouton signet de la barre d’outils, ou \kbd{Ctrl+B}. Un ruban rouge indique que la
page a un signet.

\subsection{Le panneau latéral}
Le bouton sommaire ouvre un panneau à trois onglets : \ui{Sommaire}, \ui{Signets} et \ui{Notes}.
Touchez un élément pour y aller. \ui{Tout copier en Markdown} copie tous vos surlignages et vos
notes, prêts à être collés dans un autre programme.

% ---------------------------------------------------------------------------
\section{Dictionnaire et traduction}\label{sec:lookup}

Sélectionnez un mot (jusqu’à trois) et choisissez \ui{Définir} pour voir ses définitions. Celles
de la langue du livre apparaissent en premier. Elles proviennent du Wiktionary et sont rédigées en
anglais.

Sélectionnez n’importe quel passage et choisissez \ui{Traduire}. Choisissez la langue cible dans la
liste ; l’application la retient.

\begin{note}{Connexion internet nécessaire}
Les définitions et les traductions nécessitent une connexion internet. Le texte sélectionné est
envoyé au Wiktionary ou à MyMemory ; rien d’autre n’est envoyé.
\end{note}

% ---------------------------------------------------------------------------
\section{Rechercher dans un livre}\label{sec:search}

Appuyez sur \kbd{Ctrl+F} ou sur le bouton de recherche de la barre d’outils et tapez au moins deux
lettres. La recherche ignore les majuscules et les accents. Choisissez un résultat pour y aller.
Une barre en haut affiche \emph{3 sur 27} : utilisez ses flèches, ou \kbd{F3}, pour passer au
résultat suivant. \kbd{Échap} la ferme.

% ---------------------------------------------------------------------------
\section{Affichage et accessibilité}\label{sec:display}

Ouvrez ces réglages avec le bouton de réglages de l’écran d’accueil, ou avec le bouton texte du
lecteur. Un aperçu en haut montre chaque changement immédiatement.

\shot{settings}{Affichage et accessibilité.}

\subsection{Langue}
\ui{Langue de l’application} passe l’interface en français, anglais, espagnol, allemand, italien
ou portugais instantanément. \ui{Comme l’appareil} suit la langue de votre appareil.

\subsection{Texte}
\begin{itemize}
  \item Thème : \ui{Clair}, \ui{Sépia} ou \ui{Sombre}.
  \item Police : quatre polices incluses dans l’application, pour que les livres aient le même
        aspect sur tous les appareils :
  \begin{itemize}
    \item \textbf{Literata}, conçue pour la lecture à l’écran (par défaut) ;
    \item {\garamondfont EB Garamond}, un caractère de livre classique ;
    \item {\sffamily Atkinson Hyperlegible}, créée pour que les lettres semblables se distinguent
          facilement ;
    \item {\dyslexicfont OpenDyslexic}, pensée pour les lecteurs dyslexiques.
  \end{itemize}
  \item \ui{Taille du texte}, \ui{Interligne} et \ui{Espacement des lettres}.
  \item \ui{Justifier le texte} : désactivez-le pour que les espaces entre les mots soient égaux,
        ce que beaucoup de lecteurs dyslexiques trouvent plus facile.
\end{itemize}

\subsection{Page}
\ui{Marges} et \ui{Largeur des lignes}. Les lignes plus courtes sont plus faciles à suivre.

\subsection{Écran}
\ui{Luminosité} assombrit toute l’application. \ui{Chaleur} réduit la lumière bleue, comme un mode
nuit. Les deux s’ajoutent aux réglages de votre appareil.

\ui{Rétablir les valeurs par défaut} remet tous ces réglages à leur état initial.

% ---------------------------------------------------------------------------
\section{Raccourcis clavier}

\begin{keytable}{Touches}{Action}
  \kright{} \kbd{Pg Suiv} \kbd{Espace} & Page suivante \\
  \kleft{} \kbd{Pg Préc}               & Page précédente \\
  \kup{} \kdown{}                      & Défiler un peu \\
  \kbd{Ctrl+F}                         & Rechercher dans le livre \\
  \kbd{F3}                             & Résultat suivant \\
  \kbd{Ctrl+B}                         & Ajouter ou retirer un signet sur la page \\
  \kbd{Échap}                          & Fermer la recherche ou la barre, ou revenir à la
                                         bibliothèque \\
  \kbd{Maj} + flèches                  & Étendre une sélection de texte \\
\end{keytable}

% ---------------------------------------------------------------------------
\section{Confidentialité et données}

Illuminati n’a ni compte ni suivi. Votre bibliothèque, vos notes et vos réglages sont enregistrés
uniquement sur votre appareil (sous Windows, dans \texttt{\%APPDATA\%\textbackslash Illuminati}).

L’application ne se connecte à internet que pour :
\begin{itemize}
  \item parcourir les catalogues, y chercher des livres et télécharger livres et couvertures ;
  \item définir ou traduire le texte que vous avez sélectionné, qui est envoyé au Wiktionary ou à
        MyMemory.
\end{itemize}

% ---------------------------------------------------------------------------
\section{À propos du droit d’auteur}

Tous les livres sont dans le domaine public aux États-Unis. Certains textes de Wikisource sont
plutôt sous licence libre (comme CC BY-SA), ce qui permet de les lire et de les partager librement.
D’Internet Archive, seuls les livres publiés il y a 95 ans ou plus sont affichés.

\begin{tip}{Vérifiez les lois de votre pays}
Le droit d’auteur varie d’un pays à l’autre. Un livre libre aux États-Unis peut encore être protégé
là où vous vivez ; vérifiez donc les lois de votre pays avant de télécharger.
\end{tip}

% ---------------------------------------------------------------------------
\section{Dépannage}

\begin{keytable}{Message}{Que faire}
  Impossible de joindre un catalogue & Le service est peut-être surchargé, ou votre connexion
                             coupée. Utilisez \ui{Réessayer} plus tard ; les autres catalogues
                             continuent de fonctionner. \\
  Ce titre n’a pas d’édition numérique & C’est probablement un livre audio. Cherchez une autre
                             édition. \\
  Wikisource n’a pas pu préparer ce livre & L’œuvre est trop longue pour le service de
                             Wikisource. Essayez l’un de ses volumes ou un autre catalogue. \\
  Limite quotidienne de traductions gratuites atteinte & MyMemory autorise un nombre limité de
                             traductions gratuites par jour. Réessayez demain. \\
  Le téléchargement a pris trop de temps & Vérifiez votre connexion et réessayez. \\
\end{keytable}

% ---------------------------------------------------------------------------
\section{Crédits}

Livres : Project Gutenberg, Standard Ebooks, Wikisource (via WS Export) et Internet Archive.
Définitions : Wiktionary (CC BY-SA). Traductions : MyMemory.

Polices, sous licence SIL Open Font License : Literata, EB Garamond, Atkinson Hyperlegible
(Braille Institute) et OpenDyslexic (Abbie Gonzalez). Leurs licences se trouvent dans
\ui{À propos} → \ui{Licences} dans l’application.

\end{document}
